\section{Angular momentum: an operator cross product with an explicit definition}

Orbital angular momentum is
\begin{equation}
\hat{\bm{L}}=\hat{\bm{x}}\times\hat{\bm{p}},
\end{equation}
and its components satisfy
\begin{equation}
\boxed{
\comm{\hat L_i}{\hat L_j}
=
\ii\hbar\epsilon_{ijk}\hat L_k.
}
\label{eq:Lalgebra}
\end{equation}
If one explicitly defines the ordered operator cross product by
\begin{equation}
(\hat{\bm{L}}\times\hat{\bm{L}})_i
:=
\epsilon_{ijk}\hat L_j\hat L_k,
\end{equation}
then antisymmetry of \(\epsilon_{ijk}\) removes the anticommutator part:
\begin{align}
(\hat{\bm{L}}\times\hat{\bm{L}})_i
&=
\frac{1}{2}\epsilon_{ijk}\comm{\hat L_j}{\hat L_k}
\\
&=
\ii\hbar\hat L_i.
\end{align}
Hence
\begin{equation}
\boxed{
\hat{\bm{L}}\times\hat{\bm{L}}
=
\ii\hbar\hat{\bm{L}},
}
\label{eq:Lcross}
\end{equation}
but only under the stated ordering convention. Equation~\eqref{eq:Lcross} is a consequence of the Lie algebra in Eq.~\eqref{eq:Lalgebra}, not an identity for classical gradients.
